\section{Conclusion}
\label{sec:conclusion}

In this thesis, we proposed a formulation of the multi-driver ride-sharing problem as a combination of separate single-driver scheduling problems. We modelled it as three subproblems: matching, scheduling, and conflict resolution. We formulated a three-step pipeline that first matches customers to drivers based on customer preferences, then finds each driver's schedule independently, and finally performs greedy conflict resolution to provide each driver's final conflict-free schedule. 

We found that the runtime of all three steps grows superlinearly with the number of customers. The matching and scheduling runtimes show no clear trend as the number of drivers increases, while conflict resolution becomes faster as the number of drivers increases. The scheduling step took the least time by far, at most a bit over 10 seconds, while conflict resolution took a little over 190 seconds for the scenario with 8000 customers and 100 drivers. The matching step took around 120 seconds for the largest scenarios. All in all, the largest scenarios took a bit over five minutes to solve.

Although it took a long time to solve, conflict resolution improved the objective value from the initial scheduling round by only 9--17\%. In the over-provisioned reference scenario with 1000 customers and 200 drivers, the conflict resolution step even worsened the result by almost 11\%. The solution was always optimal for the driver in question, given the current assignable set of customers, but not necessarily optimal for the entire driver fleet. Because the conflict resolution is greedy, the resulting schedules may not be optimal, and this effect can be amplified in customer-dense scenarios like the reference scenario. Also, because approximately 2/3 of customers were unserved in all but the reference scenario, all used scenarios are rather customer-dense. The order in which the driver resolves conflicts had almost no effect on the result, as the median excess over the per-scenario best ordering is below 0.11\% for all strategies. While the ascending conflicts ordering had both the smallest spread and median excess, no single ordering outperformed consistently.

As previously said, we investigated only one formulation of the ride-sharing problem. In our mathematical model, we considered only the time and coverage dimensions and modelled, for example, the distance dimension with a simple mandatory break between all rides. In reality, we would need to account for the time spent travelling between drop-off and pickup locations. Naturally, many more aspects could be investigated: traffic, travel cost, and the road network, to name a few. In the matching step, we considered only customer preferences. We could extend the preferences to include drivers' preferences, since some drivers may not want to serve certain customer types. The conflict resolution algorithm could also be improved to ensure a more globally optimal result. We could also improve the conflict-resolution metric, as our model used only pickup delay. In addition, we could improve the matching algorithm. We used Ford–Fulkerson with an Edmonds–Karp implementation, but improved algorithms already exist for solving maximum flow in a graph. \comments{A smoother ending}
